\begin{table}[htbp]
\centering
\caption{Endpoint coverage and monthly labor income}
\label{tab:latam_country_income_descriptives}
\small
\resizebox{\textwidth}{!}{%
\begin{tabular}{lrrrrrr}
\toprule
Economy & Initial & Final & Published initial & Published final & Synthetic initial & Synthetic final \\
\midrule
Argentina & 2016 & 2023 & 1,267 & 918 & 1,267 & 918 \\
Bolivia & 2016 & 2023 & 997 & 958 & 980 & 943 \\
Brazil & 2016 & 2023 & 897 & 945 & 892 & 942 \\
Chile & 2016 & 2023 & 1,174 & 1,291 & 1,160 & 1,278 \\
Colombia & 2016 & 2023 & 755 & 908 & 749 & 903 \\
Costa Rica & 2016 & 2024 & 1,328 & 1,315 & 1,298 & 1,304 \\
Dominican Republic & 2016 & 2023 & 800 & 837 & 799 & 835 \\
Ecuador & 2016 & 2024 & 843 & 708 & 838 & 694 \\
El Salvador & 2016 & 2023 & 564 & 690 & 559 & 687 \\
Guatemala & 2016 & 2022 & 614 & 634 & 609 & 629 \\
Mexico & 2016 & 2023 & 614 & 673 & 640 & 698 \\
Paraguay & 2017 & 2023 & 1,007 & 1,002 & 999 & 994 \\
Peru (Lima and Callao) & 2016 & 2023 & 927 & 829 & 925 & 825 \\
Uruguay & 2016 & 2023 & 1,130 & 1,129 & 1,129 & 1,128 \\
\bottomrule
\end{tabular}
}
\begin{minipage}{0.98\textwidth}
\footnotesize
\textit{Note:} All observations are fourth-quarter values in 2017 PPP dollars. Published income is LABLAC's reported total. Synthetic income weights the three education-group means with the published number of workers in each group.
\end{minipage}
\end{table}
